\chapter{Why it exists}

It was built as a technical challenge for a job application to Humanoid, a London company building
humanoid robots. Their brief was to compare two phone recordings of a room and identify meaningful
changes in three dimensions, while distinguishing genuine change from differences in camera
viewpoint, occlusion and incomplete observation.

That last clause is the whole exercise. Anyone can write something that notices two pictures are
not identical. The brief asks for something that knows when it is entitled to say so.

\section{Why a robot would want this}

A robot working in a warehouse, a hospital or a home builds a map of its surroundings. The world
then moves. Pallets get shifted. Doors close. Boxes appear in corridors. Furniture is rearranged
around it. The robot has to notice, without being re-shown the whole building every morning.

\begin{plainbox}
\textbf{The consequence of being wrong is not symmetrical.}

If the robot wrongly believes a chair is still there, it goes around a chair that is not there. It
wastes a few seconds.

If the robot wrongly believes a corridor is clear, because it never actually looked, it drives into
something.
\end{plainbox}

So a system like this should prefer saying \emph{I do not know} over guessing. That single
principle runs through every design decision in the project, and it is why the system is
deliberately built to under-report rather than over-report.

\section{Why it is not solved already}

Depth cameras and laser scanners make this much easier, and plenty of robots carry them. This
project deliberately uses nothing but an ordinary phone camera, because the brief said phone
recordings.

That removes the one thing that would make the problem simple, which is a direct measurement of
distance. Everything about distance here has to be inferred from how the view changes as the camera
moves, and from a model that has learned what near and far look like. Chapter~3 covers both.
